% ----------------------
\subsection{Section 3 (July 1828)}
% ----------------------
	\label{DC3}
	
	Ann Taves's book \textit{Revelatory Events} (2016, Princeton UP) has extensive discussion of this section and she tries to reconstruct a process by which JS might have received it. See my copy of her book \href{https://drive.google.com/file/d/1LfyadWG4dd4Q0FD_upWx8ULnlFvvRWKJ/view?usp=sharing}{here}; the relevant material starts on page 31, though she does talk about it before then too.
	
	% -----
	\subsubsection{Historical Background}
	% -----
		\label{DC3-HB}
		
		% --
		\paragraph{JSP}
		% --
		
			``This is the first JS revelation for which a text has survived. According to JS's history, it was obtained using the Urim and Thummim after Martin Harris lost the earliest Book of Mormon manuscript. Working in Harmony, Pennsylvania, from mid-April to mid-June 1828, JS dictated and Harris wrote what JS later called `the Book of Lehi,' translated from the gold plates. Facing ongoing opposition from his wife, Harris pressed JS to `enquire of the Lord through the Urim and Thummin' whether Harris could take the translation to his home in Palmyra, New York, to `read to his friends that peradventur he might convince them of the truth.' JS's history said that he `inquired of the Lord and the Lord said unto me that he [Harris] must not take' the manuscript. Dissatisfied with the answer, Harris made a second request, resulting in a similar inquiry and answer. Finally, `after much solicitation' from Harris, JS `again enquired of the Lord, and permission was granted him [Harris] to have the writings on certain conditions,' which included showing them only to specified family members.
			
			``With the manuscript in hand, Harris left for New York, apparently on 14 June. On 15 June, Emma Smith gave birth to a son who was either stillborn or died shortly after birth. For nearly three weeks JS cared for Emma, who in Lucy Mack Smith's words felt `to tremble upon the verge of the silent home of her infant.' As Emma began to recover in early July, she encouraged JS to `repair to Palmira, for the purpose of learning the cause of Mr. Harris absence, as well as silence.' Shortly after JS arrived at his parents' home near Palmyra, he learned that Harris had lost the manuscript. Lucy Mack Smith wrote that `sobs and groans...filled the house,' with JS `weeping and grieving like a tender infant.' He and his family soon `parted with heavy hearts,' and JS returned to Pennsylvania.
			
			``JS's history described the historical setting for the revelation: `Immediately after my return home [to Harmony] I was walking out a little distance, when Behold the former heavenly messenger appeared and handed to me the Urim and Thummin again (for it had been taken from me in consequence of my having wearied the Lord in asking for the privilege of letting Martin Harris take the writings which he lost by transgression) and I enquired of the Lord through them and obtained the folowing revelation.' It is not known when or how the text was committed to paper. Although JS may have written it himself, he dictated later revelations to scribes and may have dictated this one to either Emma or her brother Reuben Hale, both of whom served as scribes to JS during this time. The earliest extant version of the revelation...was copied into Revelation Book 1 by John Whitmer, likely in early 1831.
			
			``The revelation rebuked JS for allowing Harris to take the manuscript and commanded him to repent before he could resume the translation. According to Lucy Smith, the angel told JS that because he `had sinned in that he had delivered the manuscript into the hands of a wicked man...he would of necessity suffer the consquence's of his indiscretion.' However, after losing the plates and the spectacles, JS continued his `suplications to God without cessation' until he `had the joy and satisfaction of again receiving the urim and Thummim''' (D1:6--7).
			
		% --
		\paragraph{BSJSP}
		% --
		
			See \hyperref[EXSS-DC3-BACKGROUND]{here} for a long extract from Lucy Mack Smith's history on the circumstances surrounding this revelation and experience with Martin Harris.
		
	% -----
	\subsubsection{Versions}
	% -----
		\label{DC3-V}
	
		\begin{tabular}{p{1in}p{2in}p{1in}p{2in}}
			JSP 	& \hyperref[EMS-1828-1829-JS-JUL-1828]{July 1828}	& BC		& Chapter 2, 7--9 \\
			DC 1835 & Section 30, 156--157								& TS		& 3:14 (16 May 1842), 786 \\
			MS		& 3:6 (October 1842), 102--103						& DC 1844	& Section 30, 229--231 \\
		\end{tabular}
		
		\medskip
		
		See the \href{https://drive.google.com/file/d/1goAvNz8jS4R8slYfMG_db55Ty2z_vmgK/view?usp=sharing}{sections comparison} document for more information about the version history.

	% -----
	\subsubsection{Section 3:1} % *DC3-1 
	% -----
		\label{DC3-1}
		
			\footnote{%
				% \label{}%
				The JSP version, which comes from RB1, has this heading: ``July one Thousand Eight hundred \& Twenty Eight Given to Joseph the Seer after he had lost certan writings which he had Translated by the gift \& Power of God.'' BC has, ``1 \textit{A Revelation given to Joseph, in Harmony, Pennsylvania, July, 1828, after Martin [Harris] had lost the Manuscript of the forepart of the book of Mormon, translated from the book of Lehi, which was abridged by the hand of Mormon, saying:}''. Both the 1835 and 1844 DC insert a large heading over this section that reads, ``\textsc{On Priesthood and Calling}.'' The 1835 then has a section description: ``\textit{Revelation to Joseph Smith, jr. given July, 1828, concerning certain manuscripts on the first part of the book of Mormon, which had been taken from the possession of Martin Harris.}'' The 1844 also has a section description: ``Revelation to Joseph Smith, jr. given July, 1828, concerning certain manuscripts on the first part of the book of Mormon, which had been taken from the possession of Martin Harris.''
				}%
		THE works, and the designs, and the purposes of God cannot be frustrated, neither can they come to naught.

		\begin{framed}
		\scriptsize
			\begin{compactdesc}
				\item[JSP] Saying the \sout{words} <works> \sout{of} \& designs \& the Purposes of God cannot be frustrated neither can they come to \sout{naught} ground
				\item[BC] THE works, and the designs, and the purposes of God, can not be frustrated, neither can they come to nought,
				\item[DC 1835] 1 The works, and the designs, and the purposes of God, cannot be frustrated, neither can they come to nought,
				\item[DC 1844] 1 The works, and the designs, and the purposes of God, cannot be frustrated, neither can they come to nought,
			\end{compactdesc}
		\end{framed}

	% -----
	\subsubsection{Section 3:2} % *DC3-2 
	% -----
		\label{DC3-2}

		For God doth not walk in crooked paths, neither doth he turn to the right hand nor to the left, neither doth he vary from that which he hath said, therefore his paths are \hyperref[STRAIGHT]{straight},%
			\footnote{%
				% \label{}%
				All of these statements here suggest that the Lord takes the most efficient route to accomplish something. From our perspective, his route might look really circuitous, but that will only be because we don't see all the factors that go into his decision. Human factors, among others, affect what the most efficient route is for God.
				}
		and his course is one eternal 
			\footnote{%
				% \label{}%
				For more on the phrase ``one eternal round,'' see \hyperref[ROUND-PHRASES]{here}.
				}%
		\hyperref[ROUND]{round}.

		\begin{framed}
		\scriptsize
			\begin{compactdesc}
				\item[JSP] for God doth not wa[l]k in crooked Paths neither doth he turn to the right hand nor to the left neither doth vary from that which he hath said therefore his paths are \hyperref[STRAIT]{strait} \& his course is one eternal round
				\item[BC] for God doth not walk in crooked paths; neither doth he turn to the right hand nor to the left; neither doth he vary from that which he hath said: Therefore his paths are \hyperref[STRAIT]{strait} and his course is one eternal round.
				\item[DC 1835] for God doth not walk in crooked paths; neither doth he turn to the right hand nor to the left: neither doth he vary from that which he hath said: therefore his paths are \hyperref[STRAIT]{strait} and his course is one eternal round.
				\item[DC 1844] for God doth not walk in crooked paths; neither doth he turn to the right hand nor to the left; neither doth he vary from that which he hath said: therefore his paths are \hyperref[STRAIGHT]{straight} and his course is one eternal round.
			\end{compactdesc}
		\end{framed}

	% -----
	\subsubsection{Section 3:3} % *DC3-3 
	% -----
		\label{DC3-3}

		Remember, remember that it is not the work of God that is frustrated, but the work of men;

		\begin{framed}
		\scriptsize
			\begin{compactdesc}
				\item[JSP] Remember Remember that it is not the work of God that is frustrated but the works of men
				\item[BC] 2 Remember, remember, that it is not the work of God that is frustrated, but the work of men:
				\item[DC 1835] 2 Remember, remember, that it is not the work of God that is frustrated, but the work of men:
				\item[DC 1844] 2 Remember, remember, that it is not the work of God that is frustrated, but the work of men:
			\end{compactdesc}
		\end{framed}

	% -----
	\subsubsection{Section 3:4} % *DC3-4 
	% -----
		\label{DC3-4}

		For although a man may have many revelations, and have power to do many mighty works, yet if he boasts in his own strength, and sets at naught the counsels of God,%
			\footnote{%
				% \label{}%
				For here and below, ``set at naught'' is glossed in the 1828 dictionary as ``to slight, to disregard or despise.'' The phrase also appears at \hyperref[PROVERBS-1-25]{Proverbs 1:25}, \hyperref[1NEPHI-19-7]{1 Nephi 19:7} (twice), and \hyperref[HELAMAN-12-6]{Helaman 12:6}, besides three times in this section of the DC. The verse in Helaman is particularly helpful here: ``Notwithstanding his great goodness and his mercy towards them, they do set at naught his counsels and they will not that he should be their guide.'' People who ``set at naught'' the counsel do not want the Lord to be their guide; they want to do things on their own, or follow after their own way (Harris ``depended upon his own judgment and boasted in his own wisdom,'' which directly violates the principle given in this verse). Another very helpful example comes from this section of the DC, where the phrase appears three times (here, \hyperref[DC3-7]{verse 7}, and \hyperref[DC3-13]{verse 13}). The Lord says that Martin Harris did this, so we can look to his actions for another example. It's also possible that JS did it as well, since, according to the story, he received a ``no'' answer twice before receiving the ``yes.'' This story helps us understand why the heading ``\textsc{On Priesthood and Calling}'' was included in early versions of the DC. 
				
				When I think of ``setting at naught the counsel,'' I also think of what that guy said at BYU so many years ago, about how, if you want to be successful, you need to study righteousness.
				}
		and follows after the dictates of his own will and carnal desires,
			\footnote{%
				% \label{}%
				The JSP version has, ``he must fall to the Earth'' here.
				}%
		he must fall and incur the vengeance of a just God upon him.

		\begin{framed}
		\scriptsize
			\begin{compactdesc}
				\item[JSP] for although a man may have many Revelations \& have power to do many Mighty works yet if he boast in his own strength \& Sets at naught the councils of God \& follows after the dictates of his will \& carnal desires he must fall to the Earth \& incur the vengence of a Just God upon him
				\item[BC] for although a man may have many revelations, and have power to do many mighty works, yet, if he boasts in his own strength, and sets at nought the counsels of God, and follows after the dictates of his own will, and carnal desires, he must fall and incur the vengeance of a just God upon him.
				\item[DC 1835] for although a man may have many revelations, and have power to do many mighty works, yet, if he boasts in his own strength, and sets at nought the counsels of God, and follows after the dictates of his own will, and carnal desires, he must fall and incur the vengeance of a just God upon him.
				\item[DC 1844] for although a man may have many revelations, and have power to do many mighty works, yet, if he boasts in his own strength, and sets at nought the counsels of God, and follows after the dictates of his own will, and carnal desires, he must fall and incur the vengeance of a just God upon him.
			\end{compactdesc}
		\end{framed}

	% -----
	\subsubsection{Section 3:5} % *DC3-5 
	% -----
		\label{DC3-5}

		Behold, you have been entrusted with these things, but how strict were your commandments; and remember also the promises which were made to you,%
			\footnote{%
				% \label{}%
				The JSP version makes it seem as though the promises were about what would happen if he \textit{did} go against the commandments, rather than following the commandments.
				}
		if you did not transgress them.

		\begin{framed}
		\scriptsize
			\begin{compactdesc}
				\item[JSP] behold you have been intrusted with those things but strict was your commandment \& Remember also the Promises which were made to you if you transgressed them
				\item[BC] 3 Behold, you have been intrusted with these things, but how strict were your commandments; and remember, also, the promises which were made to you, if you did not transgress them;
				\item[DC 1835] 3 Behold, you have been intrusted with these things, but how strict were your commandments; and remember, also, the promises which were made to you, if you did not transgress them;
				\item[DC 1844] 3 Behold, you have been intrusted with these things, but how strict were your commandments; and remember, also, the promises which were made unto you, if you did not transgress them;
			\end{compactdesc}
		\end{framed}

	% -----
	\subsubsection{Section 3:6} % *DC3-6 
	% -----
		\label{DC3-6}

		And behold, how oft you have transgressed the commandments and the laws of God, and have gone on in the \hyperref[PERSUASION]{persuasions} of men.

		\begin{framed}
		\scriptsize
			\begin{compactdesc}
				\item[JSP] \& behold how oft you have transgressed the\sout{m} Laws of God \& have gone on in the Persuasions of men
				\item[BC] and behold, how oft you have transgressed the commandments and the laws of God, and have gone on in the persuasions of men:
				\item[DC 1835] and behold, how oft you have transgressed the commandments and the laws of God, and have gone on in the persuasions of men:
				\item[DC 1844] and behold, how oft you have transgressed the commandments and the laws of God, and have gone on in the persuasions of men:
			\end{compactdesc}
		\end{framed}

	% -----
	\subsubsection{Section 3:7} % *DC3-7 
	% -----
		\label{DC3-7}

		For, behold, you should not have feared man more than God. \hyperref[ALTHOUGH]{Although} men set at naught the counsels of God, and despise his words---

		\begin{framed}
		\scriptsize
			\begin{compactdesc}
				\item[JSP] for behold you should not have feared men more then God although men set at naught the councils of God \& dispise his words
				\item[BC] for behold, you should not have feared man more than God, although men set at naught the counsels of God, and despise his words,
				\item[DC 1835] for behold, you should not have feared man more than God, although men set at nought the counsels of God, and despise his words,
				\item[DC 1844] for behold, you should not have feared man more than God, although men set at nought the counsels of God, and despise his words, 
			\end{compactdesc}
		\end{framed}

	% -----
	\subsubsection{Section 3:8} % *DC3-8 
	% -----
		\label{DC3-8}

		Yet you should have been faithful; and he would have extended his arm and supported you against all the fiery darts of the adversary; and he would have been with you in every time of trouble.%
			\footnote{%
				% \label{}%
				Remember that Joseph and Emma lost a baby about a month or a month and a half before this revelation was received in 1828.
				}

		\begin{framed}
		\scriptsize
			\begin{compactdesc}
				\item[JSP] yet you should have been faithful \& he would have extended his arm \& supported you against all the firey darts of the advisary \& he would have been with you in evry time of trouble
				\item[BC] yet you should have been faithful and he would have extended his arm, and supported you against all the fiery darts of the adversary; and he would have been with you in every time of trouble. 
				\item[DC 1835] yet you should have been faithful and he would have extended his arm, and supported you against all the fiery darts of the adversary; and he would have been with you in every time of trouble.
				\item[DC 1844] yet you should have been faithful and he would have extended his arm, and supported you against all the fiery darts of the adversary; and he would have been with you in every time of trouble.
			\end{compactdesc}
		\end{framed}

	% -----
	\subsubsection{Section 3:9} % *DC3-9 
	% -----
		\label{DC3-9}

		Behold, thou art Joseph, and thou wast chosen to do the work of the Lord, but because of transgression, if thou art not aware thou wilt fall.

		\begin{framed}
		\scriptsize
			\begin{compactdesc}
				\item[JSP] behold thou art Joseph \& thou wast chosen to do the work of the Lord but because of transgression thou mayest fall
				\item[BC] 4 Behold thou art Joseph, and thou wast chosen to do the work of the Lord, but because of transgression, if thou art not aware thou wilt fall,
				\item[DC 1835] 4 Behold thou art Joseph, and thou wast chosen to do the work of the Lord, but because of transgression, if thou art not aware thou wilt fall,
				\item[DC 1844] 4 Behold thou art Joseph, and thou wast chosen to do the work of the Lord, but because of transgression, if thou art not aware thou wilt fall,
			\end{compactdesc}
		\end{framed}

	% -----
	\subsubsection{Section 3:10} % *DC3-10 
	% -----
		\label{DC3-10}

		But remember, God is merciful; therefore, repent of that which thou hast done which is contrary to the commandment which I gave you, and thou art still chosen, and art again called to the work;

		\begin{framed}
		\scriptsize
			\begin{compactdesc}
				\item[JSP] but remember God is merciful therefore repent of that which thou hast done \& he will only cause thee to be afflicted for a season \& thou art still chosen \& will \& will again be called to the work
				\item[BC] but remember God is merciful: Therefore, repent of that which thou hast done, and he will only cause thee to be afflicted for a season, and thou art still chosen, and wilt again be called to the work;
				\item[DC 1835] but remember God is merciful: therefore, repent of that which thou has done, which is contrary to the commandment which I gave you, and thou art still chosen, and art again called to the work;
				\item[DC 1844] but remember God is merciful: therefore, repent of that which thou hast done, which is contrary to the commandment which I gave you, and thou art still chosen, and art again called to the work;
			\end{compactdesc}
		\end{framed}

	% -----
	\subsubsection{Section 3:11} % *DC3-11 
	% -----
		\label{DC3-11}

		Except thou do this, thou shalt be delivered up and become as other men, and have no more gift.

		\begin{framed}
		\scriptsize
			\begin{compactdesc}
				\item[JSP] \& except Thou do this thou shalt be delivered up \& become as other men \& have no more gift
				\item[BC] and except thou do this, thou shalt be delivered up and become as other men, and have no more gift.
				\item[DC 1835] except thou do this, thou shalt be delivered up and become as other men, and have no more gift.
				\item[DC 1844] except thou do this, thou shalt be delivered up and become as other men, and have no more gift.
			\end{compactdesc}
		\end{framed}

	% -----
	\subsubsection{Section 3:12}
	% -----
		\label{DC3-12}

		And when thou deliveredst up that which God had given thee sight and power to translate,%
			\footnote{%
				\label{DC3-12-sight}%
				The manuscript in the JSP has ``when thou deliveredst up that Which that which God had given thee right to Translate'' originally. That manuscript then gets altered so the ``right'' becomes ``sight'' and then ``and power'' is added supralinearly; according to JSP-RT-FE:11, Sidney Rigdon is responsible for this change, but that same volume pages 6--7 says that JS's influence over that kind of change is unknown; in any case, JS seems to have been fine with it, because it goes into the BC and subsequent DC editions as ``sight and power.'' ``Right'' is pretty different from ``sight and power,'' and the former brings to mind a connection to \hyperref[EXOUSIA]{\grl{ἐξουσία}} and lawful governance of exercise of power. Earlier in the section, in verse 6, the Lord does tell JS that he has transgressed the ``laws of God.'' It also makes me think about how we can see JS's translation as an exercise of priesthood power in the broad sense, the one outlined by BY. And remember, the heading in the older versions of the DC for this section mentions ``priesthood.''
				}
		thou deliveredst up that which was sacred into the hands of a wicked man,

		\begin{framed}
		\scriptsize
			\begin{compactdesc}
				\item[JSP] \& when thou deliveredst up that Which that which God had given thee right to Translate thou deliveredest up that which was Sacred into the hands of a wicked man
				\item[BC] 5 And when thou deliveredst up that which God had given thee sight and power to translate, thou deliveredst up that which was sacred, into the hands of a wicked man,
				\item[DC 1835] 5 And when thou deliveredst up that which God had given thee s[i]ght and power to translate, thou deliveredst up that which was sacred, into the hands of a wicked man,
				\item[DC 1844] 5 And when thou deliveredst up that which God had given thee sight and power to translate, thou deliveredst up that which was sacred, into the hands of a wicked man,
			\end{compactdesc}
		\end{framed}

	% -----
	\subsubsection{Section 3:13} % *DC3-13 
	% -----
		\label{DC3-13}

			\footnote{%
				% \label{}%
				This verse basically says that Martin Harris did all the things listed in verse 4, that lead to a fall.
				}%
		Who has set at naught the counsels of God, and has broken the most sacred promises which were made before God, and has depended upon his own judgment and boasted in his own wisdom.

		\begin{framed}
		\scriptsize
			\begin{compactdesc}
				\item[JSP] who has Set at naught the Councils of God \& hath broken the most Sacred promises which was made before God \& hath depended upon his own Judgement \& boasted in his own \sout{arm} wisdom
				\item[BC] who has set at nought the counsels of God, and has broken the most sacred promises, which were made before God, and has depended upon his own judgment, and boasted in his own wisdom,
				\item[DC 1835] who has set at nought the counsels of God, and has broken the most sacred promises, which were made before God, and has depended upon his own judgment, and boasted in his own wisdom,
				\item[DC 1844] who has set at nought the counsels of God, and has broken the most sacred promises, which were made before God, and has depended upon his own judgment, and boasted in his own wisdom,
			\end{compactdesc}
		\end{framed}

	% -----
	\subsubsection{Section 3:14} % *DC3-14 
	% -----
		\label{DC3-14}

		And this is the reason that thou hast lost thy privileges%
			\footnote{%
				% \label{}%
				More priesthood language here, to go with ``right'' in the manuscript version of 12.
				}
		for a season---

		\begin{framed}
		\scriptsize
			\begin{compactdesc}
				\item[JSP] \& this is the reason that thou hast lost thy Privileges for a Season
				\item[BC] and this is the reason that thou hast lost thy privileges for a season,
				\item[DC 1835] and this is the reason that thou hast lost thy privileges for a season,
				\item[DC 1844] and this is the reason that thou hast lost thy privileges for a season,
			\end{compactdesc}
		\end{framed}

	% -----
	\subsubsection{Section 3:15} % *DC3-15 
	% -----
		\label{DC3-15}

		For thou hast suffered the counsel of thy director to be trampled upon from the beginning.

		\begin{framed}
		\scriptsize
			\begin{compactdesc}
				\item[JSP] for thou hast suffered that the council of thy directors to be trampeled upon from the begining
				\item[BC] for thou hast suffered the counsel of thy director to be trampled upon from the beginning.
				\item[DC 1835] for thou hast suffered the counsel of thy director to be trampled upon from the beginning.
				\item[DC 1844] for thou hast suffered the counsel of thy director to be trampled upon from the beginning.
			\end{compactdesc}
		\end{framed}

	% -----
	\subsubsection{Section 3:16} % *DC3-16 
	% -----
		\label{DC3-16}

		Nevertheless, my work shall go forth, for inasmuch as the knowledge of a Savior has come unto the world, through the testimony of the Jews, even so shall the knowledge of a Savior come unto my people---

		\begin{framed}
		\scriptsize
			\begin{compactdesc}
				\item[JSP] for as the knowledge of a Saveiour hath come to the world so shall the knowledge of my People
				\item[BC] 6 Nevertheless, my work shall go forth and accomplish my purposes, for as the knowledge of a Savior has come into the world, even so shall the knowledge of my people, 
				\item[DC 1835] 6 Nevertheless my work shall go forth, for, inasmuch as the knowledge of a Savior has come unto the world, through the testimony of the Jews, even so shall the knowledge of a Savior come unto my people; 
				\item[DC 1844] 6 Nevertheless my work shall go forth, for, inasmuch as the knowledge of a Savior has come unto the world, through the testimony of the Jews, even so shall the knowledge of a Savior come unto my people;
			\end{compactdesc}
		\end{framed}

	% -----
	\subsubsection{Section 3:17} % *DC3-17 
	% -----
		\label{DC3-17}

		And to the Nephites, and the Jacobites, and the Josephites, and the Zoramites, through the testimony of their fathers---

		\begin{framed}
		\scriptsize
			\begin{compactdesc}
				\item[JSP] the Nephities \& the Jacobites \& the Josephites
				\item[BC] the Nephites, and the Jacobites, and the Josephites, and the Zoramites, 
				\item[DC 1835] and to the Nephites, and the Jacobites, and the Josephites, and the Zoramites, through the testimony of their fathers;
				\item[DC 1844] and to the Nephites, and the Jacobites, and the Josephites, and the Zoramites, through the testimony of their fathers;
			\end{compactdesc}
		\end{framed}

	% -----
	\subsubsection{Section 3:18} % *DC3-18 
	% -----
		\label{DC3-18}

		And this testimony shall come to the knowledge of the Lamanites, and the Lemuelites, and the Ishmaelites, who dwindled in unbelief because of the iniquity of their fathers, whom the Lord has suffered to destroy their brethren the Nephites, because of their iniquities and their abominations.

		\begin{framed}
		\scriptsize
			\begin{compactdesc}
				\item[JSP] \& the Lamanites come to the \sout{Lamanites} knowledge of the Lamanites \& the Lamanites \& the Ishmaelites which dwindeled in unbelief because of the iniquities of their Fathers who hath been suffered to destroy their Brethren because of their iniquities \& their Abominations
				\item[BC] come to the knowledge of the Lamanites, and the Lemuelites and the Ishmaelites, which dwindled in unbelief, because of the iniquities of their fathers, who have been suffered to destroy their brethren, because of their iniquities, and their abominations: 
				\item[DC 1835] and this testimony shall come to the knowledge of the Lamanites, and the Lemuelites, and the Ishmaelites, who dwindled in unbelief because of the iniquity of their fathers, whom the Lord has suffered to destroy their brethren the Nephites, because of their iniquities and their abominations:
				\item[DC 1844] and this testimony shall come to the knowledge of the Lamanites, and the Lemuelites, and the Ishmaelites, who dwindled in unbelief because of the iniquity of their fathers, whom the Lord has suffered to destroy their brethren the Nephites, because of their iniquities and their abominations:
			\end{compactdesc}
		\end{framed}

	% -----
	\subsubsection{Section 3:19} % *DC3-19 
	% -----
		\label{DC3-19}

		And for this very purpose are these plates preserved, which contain these records---that the promises of the Lord might be fulfilled, which he made to his people;

		\begin{framed}
		\scriptsize
			\begin{compactdesc}
				\item[JSP] \& for this very Purpose are these Plates prepared which contain these Records that the Promises of the Lord might be fulfilled which he made to his People
				\item[BC] and for this very purpose are these plates preserved which contain these records, that the promises of the Lord might be fulfilled, which he made to his people; 
				\item[DC 1835] and for this very purpose are these plates preserved which contain these records, that the promises of the Lord might be fulfilled, which he made to his people;
				\item[DC 1844] and for this very purpose are these plates preserved which contain these records, that the promises of the Lord might be fulfilled, which he made to his people;
			\end{compactdesc}
		\end{framed}

	% -----
	\subsubsection{Section 3:20} % *DC3-20 
	% -----
		\label{DC3-20}

		And that the Lamanites might come to the knowledge of their fathers, and that they might know the promises of the Lord, and that they may believe the gospel and rely upon the \hyperref[MERIT]{merits} of Jesus Christ, and be glorified through faith in his name, and that through their repentance they might be saved. Amen.

		\begin{framed}
		\scriptsize
			\begin{compactdesc}
				\item[JSP] \& that the Lamanites might come to the knowledge of their Fathers \& that they may know the Promises of the Lord that they may believe the Gospel \& rely upon the merits of Jesus Christ \& that they might be glorified through faith in his name \& that they might repent \& be Saved Amen Received in Harmony Susquehannah
				\item[BC] and that the Lamanites might come to the knowledge of their fathers, and that they might know the promises of the Lord, and that they may believe the gospel and rely upon the merits of Jesus Christ, and be glorified through faith in his name; and that through their repentance they might be saved: Amen.
				\item[DC 1835] and that the Lamanites might come to the knowledge of their fathers, and that they might know the promises of the Lord, and that they may believe the gospel and rely upon the merits of Jesus Christ, and be glorified through faith in his name; and that through their repentance they might be saved: Amen.
				\item[DC 1844] and that the Lamanites might come to the knowledge of their fathers, and that they might know the promises of the Lord, and that they may believe the gospel and rely upon the merits of Jesus Christ, and be glorified through faith on his name; and that through their repentance they might be saved: Amen.
			\end{compactdesc}
		\end{framed}